\section{Implementation}\label{sec:implementation}

AgentChaos is an open-source Go command-line tool, released under the Apache 2.0 licence. Version 0.2.0, the version used in the pilot, supports MCP over standard input and output; an HTTP transport is planned. Figure~\ref{fig:architecture} shows its structure.

\begin{figure}[t]
\centering
\begin{tikzpicture}[
  font=\footnotesize,
  box/.style={draw, rounded corners=2pt, align=center, minimum height=7mm, inner sep=3pt},
  core/.style={box, fill=black!6},
  arr/.style={-{Stealth[length=4pt]}, thick},
  darr/.style={-{Stealth[length=4pt]}, dashed},
  lbl/.style={font=\scriptsize, midway, fill=white, inner sep=1pt}
]
  \node[box, minimum width=17mm] (agent) {Agent\\{\scriptsize (e.g.\ Claude Code)}};
  \node[box, right=6mm of agent, minimum width=24mm] (proxy) {AgentChaos proxy\\{\scriptsize switchyard core}\\{\scriptsize + fault hook}};
  \node[box, right=6mm of proxy, minimum width=17mm] (srv) {Real\\MCP server};
  \draw[arr, <->] (agent) -- (proxy);
  \draw[arr, <->] (proxy) -- (srv);

  \node[box, below=10mm of proxy, minimum width=20mm] (trace) {Trace\\(JSONL)};
  \node[box, below=8mm of trace, minimum width=26mm] (grader) {Grader\\{\scriptsize deterministic checks}};
  \node[box, left=9mm of grader, minimum width=14mm] (report) {Report};
  \node[box, above=9mm of proxy, minimum width=20mm] (yaml) {Experiment\\YAML};

  \draw[arr] (proxy) -- (trace);
  \draw[arr] (trace) -- (grader);
  \draw[arr] (grader) -- (report);
  \draw[darr] (yaml) -- node[lbl, right] {fault} (proxy);
  \draw[darr] (yaml.east) -- ($(yaml.east -| srv.east)+(3mm,0)$) |- node[lbl, pos=0.25, right] {expect} (grader.east);
\end{tikzpicture}
\caption{AgentChaos architecture. The proxy sits transparently between the agent and a real MCP server and injects the faults named in the experiment YAML. Every call is written to a trace, which the grader checks deterministically against the YAML's expectations to produce a report.}
\label{fig:architecture}
\end{figure}


\paragraph{switchyard: a shared proxy core.} The proxy is built on switchyard, a Go library that sits between an MCP client and an MCP server and, with no hooks installed, passes every message through unchanged, including capabilities, notifications and errors. A tool built on switchyard supplies \emph{hooks}. A hook sees each request before it goes upstream and its response on the way back. For a request it returns a decision: continue, continue with rewritten arguments, block with a tool error or a JSON-RPC error, answer with a synthesised result, or hang. For a response it may return a replacement. Hooks receive per-session state that persists across calls, and switchyard numbers each tool's calls in arrival order, which is the index AgentChaos uses for targeting. Hooked requests, the decisions taken for them (with any annotation a hook adds) and the responses are written as versioned JSONL events. switchyard handles the parts of a proxy that are easy to get wrong, such as message framing, matching responses to requests, and cancellation, so that a tool only states what it does to a call.

\paragraph{AgentChaos on switchyard.} The fault injector is a single switchyard hook, under three hundred lines of Go including the response rewriting. Each fault type maps to a switchyard decision: \texttt{timeout} to hang, \texttt{error} to block, \texttt{slow} to a delayed continue, and \texttt{partial} and \texttt{malformed} to a continue whose response is rewritten on the way back. Every faulted request is annotated, and the grader identifies faulted calls from these annotations in the trace. Servers with no fault get a pass-through hook so that their calls are traced too.

\paragraph{Running an agent.} The \texttt{run} command takes an experiment file and, for each run of each experiment, writes an MCP configuration that starts every server behind its own proxy, starts the agent with that configuration, waits for it to exit or for the wall-clock limit, and then merges and grades the traces. The agent is described either by a preset or by a command template with placeholders for the prompt, the MCP configuration and the allowed tools. The Claude Code preset starts the agent in headless mode with only the experiment's servers configured and only their tools pre-approved. The command writes a JSON report with per-run statuses, checks and tool-call counts, together with provenance: the AgentChaos and Go versions, the platform, and the start and end times.

\paragraph{Isolation.} Three measures keep runs independent and reduce the agent's access to the answers. First, each run has its own directory, \texttt{<suite>/<experiment>/<run>}, which is also the agent's working directory, so the experiment file and the servers' source are outside it. This is a convenience, not a sandbox: the experiment snapshot sits two levels above the run directory, and the generated MCP configuration in the run directory names its path. An agent with file access could still find the answers. In the pilot, the run directories were inside the repository, next to the server source. Claude Code's headless mode did not seal them off. The experiment servers' tools were pre-approved, and other state-changing built-in tools needed an approval that headless mode cannot give, but read-only file tools remained available. We discuss this in \S\ref{sec:threats}. Second, the agent is started in its own process group, and the whole group is killed when the run ends, so MCP servers or proxies the agent leaves behind cannot outlive the run. Third, the experiment file is copied when a suite starts and every proxy reads its fault from the copy, so an edit to the file during a suite cannot change what is injected.

\paragraph{Chaining (design, not implemented).} switchyard is intended as a shared foundation for several runtime tools in the same project, each built as a hook: IntentLock (blocking calls outside a user's capability envelope), Taintline (blocking calls whose arguments carry untrusted data), Fuse (a circuit breaker on runaway call patterns), Rewind (recording the data needed to undo calls), and Airlock (routing risky calls into a sandbox). The design composes them by chaining proxies, for example agent $\to$ Fuse $\to$ IntentLock $\to$ AgentChaos $\to$ server, with a shared chain identifier propagated through the environment so that the traces of all hops can be joined; switchyard already implements this identifier. Chaining would let an experiment check that a guard behaves correctly under a fault, for instance that a breaker trips when an agent retries a failing tool. We stress that this is a design: at the time of writing these sibling tools are specified but not built, and the pilot contains no chained experiment.
